\documentclass{article}
\usepackage{tabularx}
\begin{document}
	\begin{titlepage}
		\title{Compilers Assignment 2}
		\author{Stefano Casolari 187572}
		\maketitle
	\end{titlepage}
	\tableofcontents
	\newpage
	\section{Very Busy Expression}
	\subsection{Definition}
	An expression is considered to be \textbf{very busy} in a specific instance \textit{p} when said expression is evaluated on all path from p to the exit node without having its operands redefined along the way.
	\subsection{Data Flow Analysis}
	Backwards flow analysis proves efficient here. We need to use 
	\newline\newline
	
	\begin{math}
	In(B)=f(Out(B))\newline
	In(B)=Use(B)U(Out(B)-Def(B))
	\newline\newline
	Out(B)=\wedge In(S(B))\newline
	\wedge = \cap	
	\end{math}\newline\newline
	\begin{tabularx}{1.2\textwidth} { 
			| >{\raggedright\arraybackslash}X 
			| >{\centering\arraybackslash}X 
			| >{\raggedleft\arraybackslash}X | }
		\hline
		Domain & Set of all expressions  \\
		\hline
		Direction & Backwards\newline $In(B)=f(Out(B))$\newline $Out(B)=\wedge In(Succ(B))$  \\
		\hline
		Transfer function & $In(B)=Use(B)U(Out(B)-Def(B))$  \\
		\hline
		Meet operation & $\cap$  \\
		\hline
		Boundary Condition & $In(Exit)=\emptyset$  \\
		\hline
		Initial Interior Points & $In(B)=U$  \\
		\hline
	\end{tabularx}
	\section{Dominator Analysis}
	\subsection{Definition}
	We say that a node $B$ dominates a node $N$ if all paths from the entry to node $N$ pass through $B$. By this definition we can say that all blocks trivially dominate themselves.
	\subsection{Data Flow Analysis}
		\begin{tabularx}{1.2\textwidth} { 
			| >{\raggedright\arraybackslash}X 
			| >{\centering\arraybackslash}X 
			| >{\raggedleft\arraybackslash}X | }
		\hline
		Domain & Set of all nodes of the Graph  \\
		\hline
		Direction & Forward\newline $Out(B)=f(In(B))$\newline $In(B)=\wedge Out(Prev(B))$  \\ 
		\hline
		Transfer function & $Out(B)=Gen(B)U(In(B)-Kill(B))$  \\
		\hline
		Meet operation & $\cap$  \\
		\hline
		Boundary Condition & $Out(Entry)=\{Entry\}$  \\
		\hline
		Initial Interior Points & $Out(B)=U$  \\
		\hline
	\end{tabularx}
	\section{Constant Propagation}
\end{document}